\section{A Newton expansion for the entire harmonic difference}
\label{hn:sec:newton}

The exact-identities report \cite{RepoExact} leaves open the explicit evaluation of the
higher spectral derivatives of its entire common-shift remainder.
Its finite Stirling formula has an infinite, normally convergent
extension. This extension gives every spectral derivative, continues
the shift farther as harmonic depth increases, and identifies all shift
jets at the unshifted base point with multiple Arakawa--Kaneko functions.
Those functions are classical; the contribution here is their precise
connection with the repository's entire difference and the convergent
spectral-derivative formula. No literature-wide priority is asserted.

\subsection{Definitions and the known entire remainder}

For $\Re a>0$, use principal powers and put
\[
 \begin{gathered}
 F_a(s,u)=\sum_{n\ge0}\frac{(a+u)_n}{(a)_n(n+a)^s},\qquad
 K_a(u)=\frac{\Gamma(a)\Gamma(1+u)}{\Gamma(a+u)},\\[3pt]
 D(s,u;a)=F_a(s,u)-K_a(u)F_1(s,u).
 \end{gathered}
\]
The first series initially converges for $\Re s>1+\Re u$.
The source proves that $D$ is entire in $s$, locally near $u=0$.
If
\[
 e_r(n;a)=[u^r]\prod_{\nu=0}^{n-1}\left(1+\frac{u}{a+\nu}\right),
 \qquad M_r(s;a)=\sum_{n\ge0}e_r(n;a)(n+a)^{-s},
\]
then its depth coefficients are
\[
 D_r(s;a)=[u^r]D(s,u;a)
 =M_r(s;a)-\sum_{j=0}^r[u^j]K_a(u)\,M_{r-j}(s;1).
\]
In particular $D_0(s;a)=\zeta(s,a)-\zeta(s)$, with its removable
singularity at $s=1$ understood. Each $D_r$ is entire in $s$.

We use the source's exact Mellin remainder
\begin{align}
 R_a(t,u)
 &=-\frac{a-1}{1+u}e^{-at}
       {}_2F_1(1,a+u;2+u;1-e^{-t}),\label{hn:eq:R}\\
 D(s,u;a)&=\frac1{\Gamma(s)}\int_0^\infty
                         t^{s-1}R_a(t,u)\,dt,\qquad \Re s>0.
 \label{hn:eq:mellin}
\end{align}
To check the normalization directly, the geometric generating function is
\[
 \sum_{n\ge0}\frac{(a+u)_n}{(a)_n}e^{-(n+a)t}
 =e^{-at}{}_2F_1(1,a+u;a;e^{-t})
 =K_a(u)\frac{e^{-t}}{(1-e^{-t})^{1+u}}+R_a(t,u).
\]
The Gauss connection formula about $z=1$ gives the last equality:
the regular coefficient is $-(a-1)/(1+u)$ and the singular
coefficient is $K_a(u)$; the remaining singular Gauss factor reduces
to $e^{(a-1)t}$ before multiplication by $e^{-at}$.
This initially holds for generic parameters, then throughout the
stated germ by continuation \cite{DLMFHypergeom}.
The kernel $R_a$ is analytic at zero and decays exponentially at infinity.
Taking Mellin transforms proves \eqref{hn:eq:mellin} first in the
original common half-plane and then for $\Re s>0$ by continuation. Taylor subtraction at
zero proves the entire continuation in $s$.

Define the finite exponential polynomial
\begin{equation}\label{hn:eq:Aj}
 A_j(s)=\sum_{\ell=1}^j(-1)^{j-\ell}\binom j\ell\ell^{1-s},
 \qquad j\ge1.
\end{equation}
It is entire in $s$. At a nonpositive integer,
\begin{equation}\label{hn:eq:Stirling}
 A_j(-m)=j!\left\{\begin{matrix}m+1\\j\end{matrix}\right\},
 \qquad m\ge0,
\end{equation}
and therefore it vanishes for $j>m+1$.

\subsection{Normal convergence, including at nonpositive integers}

\begin{lemma}[Finite-difference Mellin estimate]\label{hn:lem:Aj}
For $\Re s>1-j$,
\begin{equation}\label{hn:eq:Aintegral}
 A_j(s)=\frac{(-1)^{j-1}j}{\Gamma(s)}
       \int_0^\infty t^{s-1}e^{-t}(1-e^{-t})^{j-1}\,dt.
\end{equation}
At zeros of $1/\Gamma(s)$ this is its analytic product with the
convergent integral. For every compact $S\subset\mathbb C$ there
are constants $C,P$ such that
\begin{equation}\label{hn:eq:Abound}
 \sup_{s\in S}|A_j(s)|\le C(1+\log(j+1))^P\qquad(j\ge1).
\end{equation}
Every fixed spectral derivative obeys a bound of the same form,
with constants depending on its order and $S$.
\end{lemma}

\begin{proof}
For $\Re s>0$, expand the finite binomial power and use
$j\binom{j-1}{\ell-1}=\ell\binom j\ell$.
Termwise Laplace integration gives \eqref{hn:eq:Aintegral}.
The integrand at zero is $O(t^{\Re s+j-2})$, so its integral is
holomorphic on $\Re s>1-j$. Analytic continuation gives the formula
on that larger half-plane. Thus it holds near any prescribed
nonpositive integer as soon as $j$ is sufficiently large.

Write $\sigma_-=\min_{s\in S}\Re s$ and
$\sigma_+=\max_{s\in S}\Re s$. Choose an integer $N$ such that
$\sigma_-+N-1>0$, and consider $j\ge N$. Put $q=1-e^{-1}<1$.
On $0<t\le1$,
\[
 (1-e^{-t})^{j-1}\le t^{N-1}q^{j-N},
\]
so the integral on this interval, including its factor $j$, is at
most $j q^{j-N}/(\sigma_-+N-1)$. This is exponentially small.
On $t\ge1$, put $P=\max(0,\sigma_+-1)$. The function
\[
                 p_j(t)=je^{-t}(1-e^{-t})^{j-1}
\]
is a probability density on $(0,\infty)$. With
$L_j=1+2\log(j+1)$, the integral on $[1,L_j]$ is bounded by $L_j^P$.
On $[L_j,\infty)$ use $p_j(t)\le je^{-t}$ and
\[
 j\int_{L_j}^\infty t^Pe^{-t}\,dt
 \le C_Pje^{-L_j}(1+L_j)^P
 \le C'_P(1+\log(j+1))^P.
\]
The reciprocal gamma function is bounded on $S$. The finitely many
remaining entire functions $A_j$ are harmless. Cauchy's estimate
on a slightly enlarged compact proves the derivative assertion.
\end{proof}

\begin{lemma}[Uniform binomial estimate]\label{hn:lem:binomial}
For a compact set $A\subset\mathbb C$,
\[
 \left|\binom{a-1}{j}\right|\le C_A j^{-\Re a}
 \qquad(a\in A,\ j\ge1).
\]
\end{lemma}
\begin{proof}
Use $(-1)^j\binom{a-1}{j}=\prod_{\nu=1}^j(1-a/\nu)$.
For $\nu>2\max_{a\in A}|a|$,
$\log|1-a/\nu|=-\Re a/\nu+O_A(\nu^{-2})$, uniformly.
Summing and bounding the initial finite factors proves the claim,
including when $a$ passes through a positive integer.
\end{proof}

\begin{theorem}[Normally convergent Newton expansion]\label{hn:thm:newton}
For $\Re a>0$ and $u\notin\{-1,-2,\ldots\}$,
\begin{equation}\label{hn:eq:newton}
 \boxed{\displaystyle
 D(s,u;a)=-\sum_{j=1}^\infty
                \binom{a-1}{j}\frac{A_j(s)}{u+j}.}
\end{equation}
The series converges normally on compact subsets of the indicated
domain with $s\in\mathbb C$. It extends the depth germ to a
meromorphic function of $u$, with at most simple poles at negative
integers, and
\begin{equation}\label{hn:eq:uresidue}
 \mathop{\rm Res}_{u=-j}D(s,u;a)=-\binom{a-1}{j}A_j(s).
\end{equation}
These poles are removable when the indicated residues vanish.

For every $r\ge0$,
\begin{equation}\label{hn:eq:depthnewton}
 \boxed{\displaystyle
 D_r(s;a)=(-1)^{r+1}\sum_{j=1}^\infty
          \binom{a-1}{j}\frac{A_j(s)}{j^{r+1}}.}
\end{equation}
This depth series converges normally on
\begin{equation}\label{hn:eq:shift-domain}
                  s\in\mathbb C,\qquad \Re a>-r,
\end{equation}
and supplies a jointly holomorphic continuation there.
\end{theorem}

\begin{proof}
On a compact $u$-set avoiding the negative integers,
$|u+j|^{-1}\le C/j$. The lemmas bound the summands by
$Cj^{-1-\delta}(1+\log(j+1))^P$, where
$\delta=\min\Re a>0$. In the depth series the same proof uses
$\delta=\min(\Re a+r)>0$. This proves normal convergence and
permits arbitrary fixed derivatives in $s$ and $a$.

To identify the sum, put $w=1-e^{-t}$. Euler's hypergeometric
transformation \cite{DLMFHypergeom} and then its series give
\begin{align*}
 R_a(t,u)
 &=-\frac{a-1}{1+u}e^{-t}
       {}_2F_1(2-a,1+u;2+u;w)\\
 &=e^{-t}\sum_{j=1}^\infty
        (-1)^j\binom{a-1}{j}\frac{j}{u+j}w^{j-1}.
\end{align*}
Initially $u$ is near zero and $0<w<1$. For $\Re s>0$ this
series can be integrated absolutely against $t^{s-1}\,dt$:
apply \eqref{hn:eq:Aintegral} with $s$ replaced by $\Re s$,
and the estimates already proved. Its Mellin transform is
\[
 \frac1{\Gamma(s)}\int_0^\infty t^{s-1}R_a(t,u)\,dt
 =-\sum_{j\ge1}\binom{a-1}{j}\frac{A_j(s)}{u+j}.
\]
It equals $D$ in a common domain; both sides are entire in $s$.
The normal convergence gives the continuation and residues in $u$.
Expand $(u+j)^{-1}$ on $|u|<1$ and extract depth coefficients.
The enlarged shift domain follows from direct normal convergence,
not interpolation from integer shifts.
\end{proof}

\subsection{All spectral derivatives}

\begin{corollary}[Logarithmic finite-difference formula]\label{hn:cor:jets}
For $m,k,r\ge0$,
\begin{align}
 \partial_s^kD(-m,u;a)
 &=(-1)^{k+1}\sum_{j\ge1}\frac{\binom{a-1}{j}}{u+j}
       \sum_{\ell=1}^j(-1)^{j-\ell}\binom j\ell
                     \ell^{m+1}(\log\ell)^k,
 \label{hn:eq:spectral-full}\\
 \partial_s^kD_r(-m;a)
 &=(-1)^{r+k+1}\sum_{j\ge1}\frac{\binom{a-1}{j}}{j^{r+1}}
       \sum_{\ell=1}^j(-1)^{j-\ell}\binom j\ell
                     \ell^{m+1}(\log\ell)^k.
 \label{hn:eq:spectral-depth}
\end{align}
The first formula is normally convergent for $\Re a>0$, $u$ away
from the negative integers; the second for $\Re a>-r$.
At $k=0$ use $(\log1)^0=1$.
\end{corollary}
\begin{proof}
Differentiate the normally convergent series using the compact
spectral-derivative estimates.
\end{proof}

At $k=0$ this is exactly the source's finite Stirling formula,
\[
 D(-m,u;a)=-\sum_{j=1}^{m+1}
       \left\{\begin{matrix}m+1\\j\end{matrix}\right\}
       \frac{(a-1)^{\underline j}}{u+j}.
\]
For $k\ge1$ the finite difference generally no longer vanishes at
large $j$. The answer is an absolutely convergent infinite identity;
no finite reduction to polygamma functions at arbitrary depth is
claimed. A finite value formula valid only at discrete values
of $s$ cannot itself be differentiated in $s$.

For clarity, if $[\epsilon^k]$ denotes a Laurent coefficient at
$s=-m$, the identity defining $D_r$ gives, for $k\ge0$,
\begin{align}
 [\epsilon^k]M_r(-m+\epsilon;a)
 ={}&\sum_{j=0}^r[u^j]K_a(u)\,
                 [\epsilon^k]M_{r-j}(-m+\epsilon;1)\notag\\
 &+\frac1{k!}\partial_s^kD_r(-m;a).
 \label{hn:eq:regular-reconstruction}
\end{align}
Thus the theorem determines every coefficient of the entire
shift correction. It does not by itself reduce the unshifted
height-one Laurent coefficients, or the entire outer-shift family
of the earlier report, to a finite basis of ordinary zeta derivatives.


Depth-zero special cases give useful independent checks:
\begin{align}
 \log\Gamma(a)
 &=\sum_{j\ge1}\frac{\binom{a-1}{j}}j
      \sum_{\ell=1}^j(-1)^{j-\ell}\binom j\ell\ell\log\ell,
 \label{hn:eq:loggamma}\\
 \gamma_k(a)-\gamma_k(1)
 &=-\sum_{j\ge1}\frac{\binom{a-1}{j}}j
      \sum_{\ell=1}^j(-1)^{j-\ell}\binom j\ell(\log\ell)^k.
 \label{hn:eq:Stieltjes-Newton}
\end{align}
The Stieltjes convention is
$\zeta(1+z,a)=z^{-1}+\sum_{k\ge0}(-1)^k\gamma_k(a)z^k/k!$.
The first identity follows from Lerch's formula \cite{DLMFHurwitz} and the second
from the regular expansion of $D_0$ at $s=1$. These are
specializations of a Newton/Hasse type calculus, with no claim
that the classical log-gamma or Stieltjes expansions are new.

There is a finite reduction at every nonnegative integer
\emph{depth-generating parameter} $u=q$. For $q\ge1$ and initially $\Re a>0$,
\begin{equation}\label{hn:eq:integeru}
 D(s,q;a)=\frac1{(a)_q}\sum_{\ell=1}^q
      \left[\begin{matrix}q\\\ell\end{matrix}\right]
            \{\zeta(s-\ell,a)-\zeta(s-\ell)\},
\end{equation}
where brackets denote unsigned Stirling numbers of the first kind.
Indeed $(a+q)_n/(a)_n=(n+a)_q/(a)_q$, and expand the rising
factorial polynomial for both $F_a$ and $K_aF_1$.
All spectral derivatives follow by replacing each zeta by its
order derivative. In particular
\[
 \partial_s^kD(-m,1;a)
   =\frac{\zeta^{(k)}(-m-1,a)-\zeta^{(k)}(-m-1)}a.
\]
Every zeta difference in \eqref{hn:eq:integeru} is entire,
including at its removable principal pole.

\subsection{The Arakawa--Kaneko boundary and every shift jet}

For a positive integer multi-index
$\boldsymbol{k}=(k_1,\ldots,k_p)$ use the descending-index convention
\[
 \operatorname{Li}_{\boldsymbol{k}}(z)
 =\sum_{n_1>\cdots>n_p\ge1}
       \frac{z^{n_1}}{n_1^{k_1}\cdots n_p^{k_p}}.
\]
Define the classical multiple Arakawa--Kaneko function by
\begin{equation}\label{hn:eq:AK}
 \Xi(\boldsymbol{k};s)=\frac1{\Gamma(s)}\int_0^\infty
       \frac{t^{s-1}}{e^t-1}
              \operatorname{Li}_{\boldsymbol{k}}(1-e^{-t})\,dt,
 \qquad \Re s>0.
\end{equation}
Write $\Xi_r(s)=\Xi((r);s)$ for $r\ge1$.
Arakawa--Kaneko \cite{AK1999} already define the multiple function.
They use ascending summation indices, so
$\Xi(r+1,\{1\}^{p-1};s)$ here corresponds to their
$\xi(1,\ldots,1,r+1;s)$; the vector is reversed.
At depth one there is no reversal ambiguity, and
\[
                         \Xi_1(s)=s\zeta(s+1),
\]
with a removable singularity at $s=0$.

For completeness all functions in \eqref{hn:eq:AK} are entire in $s$.
The polylogarithm is analytic at $z=0$ and begins at degree $p$,
so its kernel is analytic at $t=0$, beginning at order $p-1$.
At infinity its growth before multiplication by $(e^t-1)^{-1}$
is at most polynomial: since all indices are positive, its
positive terms are bounded by those of
$\operatorname{Li}_{1,\ldots,1}(1-e^{-t})=t^p/p!$.
Taylor subtraction at zero, followed by the zeros of $1/\Gamma(s)$,
therefore proves entireness. This is also the continuation in
the classical theory \cite{AK1999,KT2018}.

\begin{theorem}[Boundary values and mixed shift jets]\label{hn:thm:AKbridge}
For $r\ge1$, the continuation in \eqref{hn:eq:shift-domain} gives
\begin{equation}\label{hn:eq:a0}
                      D_r(s;0)=(-1)^r\Xi_r(s).
\end{equation}
For every $r\ge0$ and $p\ge1$,
\begin{equation}\label{hn:eq:a1-jets}
 \boxed{\displaystyle
 \left.\partial_a^pD_r(s;a)\right|_{a=1}
  =(-1)^{r+p}p!\,
          \Xi(r+1,\underbrace{1,\ldots,1}_{p-1};s).}
\end{equation}
The identities hold for every $s\in\mathbb C$, and arbitrary
fixed spectral derivatives can be applied to both sides.
Consequently
\begin{equation}\label{hn:eq:AK-Taylor}
 D_r(s;a)=(-1)^r\sum_{p=1}^\infty
       (1-a)^p\Xi(r+1,\{1\}^{p-1};s),
 \qquad |a-1|<r+1.
\end{equation}
The series is locally normal jointly with $s$.
\end{theorem}

\begin{proof}
For $a=0$, $\binom{-1}{j}=(-1)^j$. Substitute into the normally
convergent depth series and apply \eqref{hn:eq:Aintegral}:
\[
 D_r(s;0)=\frac{(-1)^r}{\Gamma(s)}\int_0^\infty
     t^{s-1}e^{-t}\sum_{j\ge1}\frac{(1-e^{-t})^{j-1}}{j^r}\,dt.
\]
For $\Re s>0$, termwise integration is justified absolutely.
The sum is $\operatorname{Li}_r(w)/w$, where $w=1-e^{-t}$,
giving \eqref{hn:eq:a0}. Continue in $s$.

For the shift jets let $e_j(N)$ be the elementary symmetric
polynomial of degree $j$ in $1,1/2,\ldots,1/N$, with $e_0(N)=1$
and $e_j(N)=0$ for $j>N$. The finite product
\[
 \binom{a-1}{n}=\frac{(-1)^{n-1}(a-1)}n
          \prod_{\nu=1}^{n-1}\left(1-\frac{a-1}{\nu}\right)
\]
gives
\[
 \left.\partial_a^p\binom{a-1}{n}\right|_{a=1}
      =\frac{(-1)^{n-p}p!}{n}e_{p-1}(n-1).
\]
Normal convergence permits insertion into
\eqref{hn:eq:depthnewton}. The Mellin integral then contains
\[
 \sum_{n\ge1}\frac{w^{n-1}}{n^{r+1}}e_{p-1}(n-1)
 =\frac1w\operatorname{Li}_{r+1,\{1\}^{p-1}}(w),
\]
with overall sign $(-1)^{r+p}p!$. Absolute integration follows
from positive coefficients and the logarithmic growth
$e_{p-1}(n-1)=O((1+\log n)^{p-1})$.
This proves \eqref{hn:eq:a1-jets} on $\Re s>0$, then everywhere.

Finally $D_r(s;1)=0$, and $D_r$ is holomorphic for $\Re a>-r$.
Its Taylor series centered at $1$ therefore converges throughout
$|a-1|<r+1$. Cauchy estimates on compact $s$-sets prove joint
local normality.
\end{proof}

The depth-zero case of \eqref{hn:eq:a1-jets} reduces to
\[
 \Xi(\underbrace{1,\ldots,1}_{p};s)
        =\frac{(s)_p}{p!}\zeta(s+p).
\]
The boundary formula \eqref{hn:eq:a0}, however, excludes $r=0$.
The original domain does not include $a=0$, and
\[
 D_0(s;a)=a^{-s}+\zeta(s,a+1)-\zeta(s)
\]
has no jointly holomorphic extension there for general $s$.
A fictitious $\Xi_0$ with a nondecaying Mellin kernel cannot
remove this obstruction.

Some exact examples, with $\gamma_k=\gamma_k(1)$, are
\begin{align*}
 D_1(s;0)&=-s\zeta(s+1),\\
 D_1(0;0)&=-1,&
 \partial_sD_1(0;0)&=-\gamma,&
 \partial_s^2D_1(0;0)&=2\gamma_1,\\
 \partial_s^kD_1(0;0)&=(-1)^k k\gamma_{k-1}\quad(k\ge1),\\
 \partial_sD_1(-1;0)&=\frac12-\frac12\log(2\pi),\\
 \partial_sD_1(-2;0)&=\frac1{12}+2\zeta'(-1),\\
 \partial_aD_1(s;1)&=\Xi_2(s),&
 \partial_a^2D_1(s;1)&=-2\Xi(2,1;s).
\end{align*}
They are identities for the continued remainder, rather than
ordinary sums of its divergent defining series.

More generally, for an integer $N\ge0$ and depth $r\ge N+1$,
\begin{equation}\label{hn:eq:negative-shifts}
 D_r(s;-N)=\frac{(-1)^r}{N!}\sum_{h=0}^N
     \left[\begin{matrix}N+1\\h+1\end{matrix}\right]\Xi_{r-h}(s).
\end{equation}
Indeed $\binom{-N-1}{j}=(-1)^j\binom{N+j}{N}$ and
\[
 \binom{N+j}{N}=\frac1{N!}\prod_{\nu=1}^N(j+\nu)
   =\frac1{N!}\sum_{h=0}^N
       \left[\begin{matrix}N+1\\h+1\end{matrix}\right]j^h.
\]
Insert this polynomial into the convergent Mellin series. The
condition $r>N$ makes all polylogarithm indices positive and
places the shift inside $\Re a>-r$. For example
\[
                   D_2(s;-1)=\Xi_2(s)+s\zeta(s+1).
\]

Specializing the convergent Taylor expansion to $a=0$ also gives
$\Xi_r(s)=\sum_{p\ge1}\Xi(r+1,\{1\}^{p-1};s)$ for $r\ge1$,
normally on compact $s$-sets. Its elementary polylogarithmic
explanation is $\sum_{p\ge1}e_{p-1}(n-1)=n$.
This is recorded as a consequence of the shift identification,
with no new-priority assertion for the height-one summation itself.

\subsection{A boundary corollary for generalized Stieltjes constants}

The two-shift bridge in the source admits the same proof. Write
$v=b-a$, $c=1+v$, and initially assume $\Re a,\Re b,\Re c>0$.
Define
\[
 D_{a,b}(s,u)=\sum_{n\ge0}\frac{(a+u)_n}{(a)_n(n+b)^s}
   -K_a(u)\sum_{n\ge0}\frac{(1+u)_n}{n!(n+c)^s},
 \qquad D_{a,b,r}(s)=[u^r]D_{a,b}(s,u).
\]
Its entire Mellin remainder is $e^{-vt}R_a(t,u)$.
Put
\[
 A_j(s;v)=\sum_{\ell=1}^j(-1)^{j-\ell}\binom j\ell
                     \ell(\ell+v)^{-s},\qquad \Re v>-1.
\]
The Laplace calculation gives
\[
 A_j(s;v)=\frac{(-1)^{j-1}j}{\Gamma(s)}
       \int_0^\infty t^{s-1}e^{-(1+v)t}(1-e^{-t})^{j-1}\,dt,
 \qquad \Re s>1-j.
\]
Then
\begin{equation}\label{hn:eq:twodepth}
 D_{a,b,r}(s)=(-1)^{r+1}\sum_{j\ge1}\binom{a-1}{j}
                           \frac{A_j(s;b-a)}{j^{r+1}},
\end{equation}
normally on the connected domain
\begin{equation}\label{hn:eq:two-domain}
              \Re(1+b-a)>0,\qquad \Re b>-r,\qquad s\in\mathbb C.
\end{equation}
This supplies a continuation without a restriction on $\Re a$.

Here is the convergence detail needed for this enlarged domain.
On compact $s,v$ sets with $\Re v>-1$, and $0<\epsilon<\min\Re(1+v)$ sufficiently small, the modified Mellin estimate is
$|A_j(s;v)|\le C_\epsilon j^{-\Re v+\epsilon}$.
On $(0,1)$ the previous proof is unchanged; on $t\ge1$ absorb
the necessary power of $t$ into $C_\epsilon e^{\epsilon t}$ and use
\[
 j\int_0^\infty e^{-(1+\Re v-\epsilon)t}
                  (1-e^{-t})^{j-1}\,dt
 =jB(1+\Re v-\epsilon,j)
 \le C_\epsilon j^{-\Re v+\epsilon}.
\]
The last bound follows directly by bounding
$(1-x)^{j-1}$ by $e^{-(j-1)x}$ in the beta integral.
Since $a+v=b$, choose $\epsilon$ below half of
$\min(\Re b+r)$ on the compact under consideration. The series
is then bounded by a summable power
$Cj^{-1-(\Re b+r)+\epsilon}$. Identification uses the same
kernel expansion as before. All fixed spectral derivatives are
allowed by normal convergence.

For $r\ge1$, $\Re x>0$, define the classical Hurwitz extension
\[
 \Xi_r(s,x)=\frac1{\Gamma(s)}\int_0^\infty
       \frac{e^{-xt}}{1-e^{-t}}t^{s-1}
                        \operatorname{Li}_r(1-e^{-t})\,dt.
\]
This convention agrees with Coppo--Candelpergher \cite{CC2010}, and
$\Xi_r(s,1)=\Xi_r(s)$. It is entire in $s$ by the same Taylor
subtraction argument. The point $a=0$, $b=x-1$ lies in
\eqref{hn:eq:two-domain} for $r\ge1$, so
\begin{equation}\label{hn:eq:generalized-boundary}
                  D_{0,x-1,r}(s)=(-1)^r\Xi_r(s,x).
\end{equation}
At depth one, $\operatorname{Li}_1(1-e^{-t})=t$ gives
\begin{equation}\label{hn:eq:Hurwitz-boundary}
 D_{0,x-1,1}(s)=-s\zeta(s+1,x)=\partial_x\zeta(s,x).
\end{equation}
The right side is entire in $s$, including at zero.

\begin{corollary}[Generalized Stieltjes boundary jets]\label{hn:cor:Stieltjes}
For $\Re x>0$ and $k\ge1$,
\begin{equation}\label{hn:eq:generalized-Stieltjes}
 \boxed{\displaystyle
 \left.\partial_s^kD_{0,x-1,1}(s)\right|_{s=0}
      =(-1)^k k\,\gamma_{k-1}(x).}
\end{equation}
The zeroth derivative is $-1$. For $m\ge1$,
\[
 \left.\partial_s^kD_{0,x-1,1}(s)\right|_{s=-m}
 =m\zeta^{(k)}(1-m,x)-k\zeta^{(k-1)}(1-m,x),
\]
with the last term omitted when $k=0$.
At $k=1,m=0$ the boundary function is $\psi(x)$, and its
antiderivative is $\log\Gamma(x)$.
\end{corollary}
\begin{proof}
Insert the defining Laurent expansion of the Hurwitz zeta
function in \eqref{hn:eq:Hurwitz-boundary}. The factor $s$
removes its pole before coefficient extraction.
\end{proof}

Every one of these boundary jets also has a finite repeated
antiderivative. For $m,k\ge0$ and $q\ge1$, define coefficients by
\[
 \prod_{\nu=1}^{q-1}\left(1-\frac z{m+\nu}\right)^{-1}
     =\sum_{h\ge0}h_h^{[m,q-1]}z^h
\]
(the product is $1$ at $q=1$). Then
\begin{equation}\label{hn:eq:boundary-primitives}
 \mathcal P_{m,k}^{(q)}(x)=\frac{k!}{(m+1)_{q-1}}
       \sum_{j=0}^k\frac{h_{k-j}^{[m,q-1]}}{j!}
                    \zeta^{(j)}(1-m-q,x)
\end{equation}
satisfies
\[
 \frac{d^q}{dx^q}\mathcal P_{m,k}^{(q)}(x)
       =\left.\partial_s^kD_{0,x-1,1}(s)\right|_{s=-m}.
\]
Indeed a $q$-fold primitive of $-s\zeta(s+1,x)$ is
$\zeta(s-q+1,x)/(1-s)_{q-1}$. Its denominator is nonzero near
$s=-m$, and differentiating $k$ times in $s$ proves the formula.
The coefficients are complete homogeneous symmetric polynomials
in $1/(m+1),\ldots,1/(m+q-1)$, equivalently Bell polynomials in
generalized harmonic-number differences. Other $q$-fold
primitives differ by a polynomial of degree less than $q$.

\subsection{Scope and further questions}

No substantive error was found in the inspected fully-shifted
and outer-shift harmonic arguments. Their restriction of the
finite coefficient rule to the principal part and finite part
is correct. The new formulas identify the missing analytic
information in higher spectral derivatives. Three extension
requirements are explicit: a discrete finite value formula
cannot be differentiated in $s$; depth and spectral manipulations
must use the already entire remainder instead of disregarding
the moving polar divisors of $F_a$; and the common-shift $a=0$
Arakawa--Kaneko formula requires positive depth.

Further questions include identifying finite combinations of
spectral jets that reduce to ordinary Hurwitz-zeta derivatives;
determining the exact shift singularities on $\Re a=-r$;
developing stable evaluation of the finite differences, since
naive binomial summation loses precision; comparing the present
shift deformation systematically with existing generalized
Arakawa--Kaneko families; and seeking multiple-zeta or
iterated-integral identities for mixed spectral and shift
derivatives beyond the standard integer specializations.

